% Generated draft: requires mathematical and source review.
\documentclass[11pt]{article}
\usepackage[margin=25mm]{geometry}
\usepackage{fontspec}
\setmainfont{Latin Modern Roman}
\usepackage{amsmath,amssymb,mathtools,unicode-math}
\setmathfont{Latin Modern Math}
\usepackage{enumitem,needspace,xurl,hyperref}
\setlength{\parindent}{0pt}
\setlength{\parskip}{5pt}
\setlength{\emergencystretch}{3em}
\newcommand{\N}{\mathbb{N}}
\newcommand{\Z}{\mathbb{Z}}
\newcommand{\Q}{\mathbb{Q}}
\newcommand{\R}{\mathbb{R}}
\newcommand{\C}{\mathbb{C}}
\newcommand{\F}{\mathbb{F}}
\newcommand{\A}{\mathbb{A}}
\newcommand{\PP}{\mathbb{P}}
\newcommand{\E}{\mathbb{E}}
\newcommand{\eps}{\varepsilon}
\newcommand{\dd}{\,\mathrm{d}}
\begin{document}
\section*{UnsolvedMath ID: 700436}
\textbf{Source problem code:} catalogue entry
\subsection{Definitions and mathematical statement}
Let \(p\) be a prime and let \(\sigma\in\Z\) satisfy \(p\nmid\sigma\). Interpret the exponent \(2/\sigma\) in the supplied presentation as division modulo \(p\): choose \(r\in\Z\) with \(\sigma r\equiv2\pmod p\), and set
\[
G_{p,\sigma}=\left\langle a,b\ \middle|\ a^p=(ab)^3=b^2=(a^\sigma b a^r b)^2=1\right\rangle.
\]
The relation \(a^p=1\) makes this presentation independent of the representative chosen for \(r\), and it also makes \(\sigma\) depend only on its residue modulo \(p\). This modular interpretation of \(2/\sigma\) is needed to read the presentation as an ordinary group presentation; the source does not explicitly define that notation.

Write \(\operatorname{PSL}_2(p)=\operatorname{SL}_2(\F_p)/Z(\operatorname{SL}_2(\F_p))\). The source asserts the existence of an epimorphism \(\pi_{p,\sigma}:G_{p,\sigma}\twoheadrightarrow\operatorname{PSL}_2(p)\). For the epimorphism intended by the source, put
\[
N_{p,\sigma}=\ker\pi_{p,\sigma},\qquad
N_{p,\sigma}^{\mathrm{ab}}=N_{p,\sigma}/[N_{p,\sigma},N_{p,\sigma}],
\]
where \([N,N]\) is the subgroup generated by the commutators \(x^{-1}y^{-1}xy\) with \(x,y\in N\). Thus the asserted sequence is
\[
1\longrightarrow N_{p,\sigma}\longrightarrow G_{p,\sigma}
\xrightarrow{\pi_{p,\sigma}}\operatorname{PSL}_2(p)\longrightarrow1.
\]
The epimorphism is not specified in the supplied record, nor is independence of the resulting abelianization from its choice established there. The following questions therefore refer to the intended quotient map.

\begin{enumerate}
\item[(a)] Is the set
\[
\left\{p\text{ prime}:\ \text{for every }\sigma\in\Z\text{ with }p\nmid\sigma,\quad |N_{p,\sigma}^{\mathrm{ab}}|<\infty\right\}
\]
infinite?
\item[(b)] Is there an arithmetic condition on \(\sigma\) ensuring that \(N_{p,\sigma}^{\mathrm{ab}}\) is finite? The source does not specify the intended scope of such a condition, its permitted dependence on \(p\), or what additional information it should provide beyond the sufficient example below.
\end{enumerate}

As supplied background, the record states that for every prime \(p>2\) there is an admissible \(\sigma\) with \(N_{p,\sigma}=1\), and gives \(\sigma=4\) as an example. It also states that for some primes an admissible \(\sigma\) gives infinite abelianization, citing \(p=41\) with \(\sigma^2\equiv2\pmod{41}\), whereas for some other primes the abelianization is finite for every admissible \(\sigma\), citing \(p=43\). The wording of the \(p=41\) example does not explicitly distinguish existence of a suitable root from a claim about both roots. These background assertions are retained as supplied, without independent verification.
\subsection{Short English statement}
For Mennicke's presented groups and their intended quotients onto \(\operatorname{PSL}_2(p)\), ask whether infinitely many primes give kernels with finite abelianization for every admissible parameter, and seek an arithmetic condition guaranteeing this finiteness. The modular exponent notation, choice of quotient map, and intended scope of the second question require clarification.
\subsection{Sources}
\begin{itemize}
\item Supplied local UnsolvedMath registry record: permanent ID \(700436\), exact source code \texttt{catalogue entry}, title \emph{Kourovka Notebook Problem 7.39}; upstream revision \texttt{d371ee917ac9647fbb06cd6fc529deec6f3dda82}. The registry labels the problem open; this is supplied metadata, not an independently verified status assessment.
\item Attributed in the supplied record to J. Mennicke, proposed in 1980. Source: \emph{The Kourovka Notebook (Unsolved Problems in Group Theory)}, 7th issue (1980), Problem 7.39, as recorded in the 21st issue edition (2026), \texttt{arXiv:1401.0300v48}. The supplied collection description names E. I. Khukhro and V. D. Mazurov as editors. This reference and the accompanying background assertions have not been independently checked.
\end{itemize}
\end{document}
